% =============================================================================
% Sección II — Representación de Documentos
% =============================================================================
\section{Representación de Documentos}
% Responsable: Seba

\subsection{Corpus utilizado}

El punto de partida del sistema es el corpus de texto. Para este proyecto se seleccionaron \textbf{12 documentos breves} en español, cada uno cubriendo un tema diferente dentro de las Ciencias de la Computación. Los documentos se almacenan en \texttt{data/documentos/corpus.txt}, con el formato de una oración por línea, lo que simplifica su carga y permite que \texttt{cargar\_documentos()} los trate como 12 entidades independientes.

\begin{table}[H]
\centering
\caption{Corpus de 12 documentos seleccionados para el proyecto.}
\label{tab:corpus}
\begin{tabular}{|c|l|}
\hline
\textbf{ID} & \textbf{Tema del documento} \\
\hline
Doc 1  & Inteligencia Artificial \\
Doc 2  & Redes Neuronales Artificiales \\
Doc 3  & Aprendizaje Automático (Machine Learning) \\
Doc 4  & Ciberseguridad \\
Doc 5  & Sistemas Operativos \\
Doc 6  & Computación en la Nube \\
Doc 7  & Desarrollo de Software \\
Doc 8  & Bases de Datos Relacionales \\
Doc 9  & Programación Orientada a Objetos \\
Doc 10 & Algoritmos de Búsqueda y Ordenamiento \\
Doc 11 & Procesamiento del Lenguaje Natural \\
Doc 12 & Robótica \\
\hline
\end{tabular}
\end{table}

La diversidad temática del corpus es intencional: al cubrir áreas distintas de la computación, los vectores de los documentos deberían ser en su mayoría ortogonales entre sí, salvo entre documentos relacionados (como Doc~1, Doc~3 y Doc~11, que comparten el ámbito de la Inteligencia Artificial). Esto permite evaluar la capacidad del buscador para discriminar entre temas.

\subsection{Pipeline de preprocesamiento}

Antes de vectorizar, cada texto bruto pasa por un pipeline de limpieza implementado en \texttt{src/preprocesamiento.py}. Este pipeline es fundamental: sin él, palabras como ``Inteligencia'' e ``inteligencia'' serían tratadas como términos distintos, inflando el vocabulario innecesariamente.

\begin{enumerate}
    \item \textbf{Normalización a minúsculas} (\texttt{limpiar\_texto}): todas las letras se convierten a minúscula. Esto garantiza que ``Datos'', ``DATOS'' y ``datos'' se traten como el mismo término.
    
    \item \textbf{Limpieza de caracteres} (\texttt{limpiar\_texto}): se eliminan signos de puntuación, números y caracteres especiales mediante la expresión regular \texttt{[\^{}a-záéíóúüñ\textbackslash{}s]}, que conserva únicamente letras del alfabeto español (incluyendo vocales acentuadas y la ñ) y espacios en blanco.
    
    \item \textbf{Tokenización} (\texttt{tokenizar}): el texto limpio se divide en palabras individuales usando \texttt{str.split()}, que separa por espacios en blanco. Esta es la opción adecuada para el modelo Bag-of-Words, donde no se requiere análisis morfológico.
    
    \item \textbf{Remoción de stopwords} (\texttt{remover\_stopwords}): se eliminan las palabras vacías (artículos, preposiciones, conjunciones) que no aportan contenido semántico. Se utiliza la constante \texttt{PALABRAS\_VACIAS\_ES}, un conjunto (\textit{set}) de 75 palabras funcionales del español definido en el propio módulo. El uso de \textit{set} en vez de lista reduce la búsqueda de $O(n)$ a $O(1)$, lo cual es relevante cuando se procesan miles de tokens.
\end{enumerate}

Ejemplo del pipeline aplicado al primer documento:

\begin{itemize}
    \item \textbf{Texto original:} ``La inteligencia artificial es una rama de la informática que busca crear sistemas capaces de realizar tareas que normalmente requieren inteligencia humana...''
    \item \textbf{Tokens resultantes:} \texttt{[`inteligencia', `artificial', `rama', `informática', `busca', `crear', `sistemas', `capaces', `realizar', `tareas', `normalmente', `requieren', `inteligencia', `humana', `reconocimiento', `voz', `toma', `decisiones']}
\end{itemize}

El preprocesamiento redujo el documento de $\approx$30 palabras a 18 tokens útiles, eliminando artículos y preposiciones que solo añadirían ruido al vector.

\subsection{Construcción de la Matriz Documento-Término}

\subsubsection{Vocabulario generado}

La función \texttt{construir\_vocabulario()} recorre los 12 documentos preprocesados y extrae todos los términos únicos, ordenándolos alfabéticamente. Este ordenamiento no es arbitrario: garantiza que ante el mismo corpus, el vocabulario generado sea siempre idéntico, haciendo la representación vectorial \textbf{reproducible y determinista}.

El corpus completo, tras el preprocesamiento, produjo un vocabulario de \textbf{154 términos únicos}, con un total de 193 tokens (contando repeticiones). Esto significa que, en promedio, cada documento contiene $193 / 12 \approx 16$ tokens útiles.

\subsubsection{Construcción de la matriz}

La función \texttt{crear\_matriz\_documento\_termino()} construye la matriz $M \in \mathbb{R}^{12 \times 154}$ donde:
\[
M_{ij} = \text{número de veces que aparece el término } t_j \text{ en el documento } d_i
\]

La implementación parte de una matriz de ceros de tipo \texttt{numpy.float64} y recorre cada documento una sola vez para incrementar el contador de cada término encontrado. El uso de \texttt{float64} (en lugar de enteros) es una decisión de diseño deliberada: permite que la misma matriz soporte en el futuro esquemas de ponderación como TF-IDF, que producen valores decimales, sin necesidad de cambiar la estructura de datos.

La Tabla~\ref{tab:matriz_ejemplo} muestra un subconjunto reducido de la matriz real (4 documentos × 6 términos seleccionados) para ilustrar la estructura:

\begin{table}[H]
\centering
\caption{Submatriz ejemplo: 4 documentos × 6 términos representativos.}
\label{tab:matriz_ejemplo}
\begin{tabular}{|l|c|c|c|c|c|c|}
\hline
\textbf{Documento} & \texttt{artificial} & \texttt{aprendizaje} & \texttt{bases} & \texttt{datos} & \texttt{inteligencia} & \texttt{sistemas} \\
\hline
Doc 1 (IA)            & 1 & 0 & 0 & 0 & 2 & 1 \\
Doc 3 (ML)            & 1 & 1 & 0 & 0 & 1 & 1 \\
Doc 8 (BD)            & 0 & 0 & 1 & 2 & 0 & 0 \\
Doc 12 (Robótica)     & 0 & 0 & 0 & 0 & 0 & 0 \\
\hline
\end{tabular}
\end{table}

Obsérvese cómo Doc~1 y Doc~3 comparten los términos ``artificial'', ``inteligencia'' y ``sistemas'', mientras que Doc~8 se orienta exclusivamente hacia ``bases'' y ``datos''. Doc~12, al hablar de robótica, no comparte ninguno de estos 6 términos, por lo que toda su fila es cero en esta submatriz.

\subsection{Sparsidad de la Matriz Documento-Término}

Una característica estructural inevitable del modelo Bag-of-Words sobre corpus de texto real es la \textbf{sparsidad}: la gran mayoría de las celdas de la matriz tendrán valor 0, porque cada documento utiliza solo una fracción pequeña del vocabulario total.

Formalmente, la sparsidad se define como:
\[
\text{sparsidad} = \frac{\text{número de celdas con valor } 0}{m \times n} \times 100\%
\]

En nuestra implementación:
\[
\text{sparsidad} = \frac{12 \times 154 - 193}{12 \times 154} \times 100\% = \frac{1{,}655}{1{,}848} \times 100\% \approx \mathbf{89.8\%}
\]

Esto significa que \textbf{casi 9 de cada 10 celdas} de la Matriz Documento-Término contienen un 0, lo cual es completamente esperable y correcto:

\begin{itemize}
    \item Cada documento breve usa aproximadamente 16 términos únicos de un vocabulario de 154.
    \item Un documento sobre Robótica no menciona términos de Bases de Datos, y viceversa.
    \item La sparsidad es una consecuencia directa de la \textbf{especificidad temática} de los documentos.
\end{itemize}

Desde el punto de vista computacional, la sparsidad tiene implicaciones importantes: si el corpus creciera a miles de documentos con un vocabulario de decenas de miles de términos, almacenar la matriz densa en memoria sería ineficiente. Para esos casos, existen representaciones compactas (matrices esparsas de \texttt{scipy.sparse}) que solo almacenan los valores no nulos, reduciendo el uso de memoria en proporción a la sparsidad. En nuestro caso, con 12 documentos y 154 términos, la matriz densa de \texttt{numpy} es perfectamente adecuada.
